\chapter{Packaging, release CI, and how a bundle is proven}\label{ch:packaging}

Four artifacts ship: a Windows installer, a Linux \texttt{.deb}, a tarball of the frozen bundle, and a
wheel. The first three come from a PyInstaller bundle; the wheel does not contain Qt at all --- it
declares the dependency and pip fetches it. Understanding which artifact carries what is the point of
this chapter, because a bug that only shows in one of them is invisible until a Station installs it.

\section{The bundle}

\texttt{packaging/\allowbreak{}build\_\allowbreak{}pyinstaller.py} builds the frozen application. It makes a throwaway venv,
installs into it, runs \texttt{packaging/\allowbreak{}pyinstaller/\allowbreak{}ate.spec}, and can smoke-test the result:

\begin{lstlisting}[language=ATEShell]
python packaging/build_pyinstaller.py --smoke
python packaging/build_pyinstaller.py --verify dist/ate   # check an existing bundle: no build, no run
\end{lstlisting}

\texttt{verify\_bundle()} is the gate a bundle must pass: it refuses a bundle that does not carry the
licence files (\texttt{\_internal/licences}, four files generated at build time by
\texttt{packaging/\allowbreak{}licence/\allowbreak{}notices.py --out build/\allowbreak{}licences}), the report template, the bundled Lua
modules, the WebEngine binary and its \texttt{.pak} and \texttt{icudtl.dat} resources --- the gate
checks that some are there, not how many. The excludes list in \texttt{ate.spec}
has its own history: PySide6's \texttt{QtWebEngineWidgets} imports \texttt{PySide6.QtWebChannel} on its
own, and keeping it excluded produced a bundle that looked fine and then died in
\texttt{ate \_pdf} with ``could not import module 'PySide6.QtWebChannel'''. It is out of the list, and
\texttt{test\_\allowbreak{}the\_\allowbreak{}spec\_\allowbreak{}does\_\allowbreak{}not\_\allowbreak{}exclude\_\allowbreak{}what\_\allowbreak{}webengine\_\allowbreak{}imports\_\allowbreak{}itself} runs the WebEngine import and
compares what it pulls in with the excludes list --- so the list cannot drift again without failing.

The real build is gated: \texttt{ATE\_\allowbreak{}BUILD\_\allowbreak{}TESTS=1 pytest -m build} runs
\texttt{tests/\allowbreak{}test\_\allowbreak{}pyinstaller\_\allowbreak{}build\_\allowbreak27a.py}, which takes minutes. It is not run by the default suite,
and the manual says so rather than implying it is.

\section{Linux: the \texttt{.deb}, the tarball, the wheel}

\texttt{packaging/\allowbreak{}linux/\allowbreak{}build\_\allowbreak{}linux.py --build --dist dist --out dist} produces the \texttt{.deb} and the
tarball from a built bundle. The \texttt{.deb} declares \texttt{DEPENDS}, the list of system libraries the
bundle cannot carry itself; \texttt{WHEEL\_SYSTEM\_LIBS} is the list for the wheel, which ships no Qt and
therefore needs the system to provide the xcb and GL libraries QtWebEngine links against. Six xcb
libraries (cursor, icccm, keysyms, shape, xkb, xkbcommon-x11) were added to that list after a bare
\texttt{debian:stable} container aborted with ``Could not load the Qt platform plugin "xcb"''.

\texttt{packaging/\allowbreak{}linux/\allowbreak{}install.sh --prefix /\allowbreak{}opt/\allowbreak{}ate --bin /\allowbreak{}usr/\allowbreak{}bin/\allowbreak{}ate} is the tarball's installer;
\texttt{packaging/\allowbreak{}linux/\allowbreak{}deb/\allowbreak{}postinst} and \texttt{postrm} are the \texttt{.deb}'s hooks (the \texttt{.deb}
is idempotent about the \texttt{ate} group; \texttt{apt purge} removes the config file).

The wheel is built with \texttt{build}. \texttt{--no-isolation} uses the interpreter already present
instead of making a throwaway venv to download setuptools into, which is what a machine with no
network to PyPI needs:

\begin{lstlisting}[language=ATEShell]
python -m build --wheel --no-isolation -o dist
\end{lstlisting}

Its version comes from the git tag, computed once by \texttt{src/ate/\_version.py} so build time and run
time agree; untagged builds fall back to \texttt{0.0.0-dev} and a build never fails because git is
unavailable. Package data matters: a bare-key bug in \texttt{pyproject.toml} meant \texttt{no wheel had
built since packet 06b}, and \texttt{tests/\allowbreak{}test\_\allowbreak{}wheel\_\allowbreak{}build\_\allowbreak28a.py} now reads the parsed config rather
than the source text, so a non-\texttt{.py} file the wheel must ship cannot go missing silently.

\section{Windows}

\texttt{packaging/\allowbreak{}windows/\allowbreak{}ate.iss} is the Inno Setup script;
\texttt{packaging/\allowbreak{}windows/\allowbreak{}build\_\allowbreak{}installer.py} compiles it. On Linux the only thing you can do is see the
command:

\begin{lstlisting}[language=ATEShell]
python packaging/windows/build_installer.py --dry-run --dist dist
\end{lstlisting}

which prints the exact \texttt{iscc} line the CI will run. \texttt{tests/\allowbreak{}test\_\allowbreak{}inno\_\allowbreak{}installer\_\allowbreak27b.py}
parses \texttt{ate.iss} and asserts against that command, and
\texttt{tests/\allowbreak{}test\_\allowbreak{}release\_\allowbreak{}ci\_\allowbreak28c.py} checks that the flags the CI passes to the installer
(\texttt{packaging/\allowbreak{}windows/\allowbreak{}smoke\_\allowbreak{}installer.py}'s \texttt{INSTALL\_FLAGS}) are the ones the manuals tell an
administrator to use.

\section{The manuals, and where they end up}

Three manuals are built before anything is packaged: the test writer's, the station administrator's
and this one. \texttt{docs/\allowbreak{}manuals/\allowbreak{}build.py} compiles each with \texttt{latexmk} in a TeX Live
image pinned by digest --- \texttt{latest} moves, and a release has to be rebuildable --- stamps the
version being released on the title page (the string \texttt{ate --version} prints, from
\texttt{src/\allowbreak{}ate/\allowbreak{}\_\allowbreak{}version.py}), and writes the PDFs into
\texttt{src/\allowbreak{}ate/\allowbreak{}manuals}, which \texttt{.gitignore} keeps out of the repository. That one staged
directory is what makes them ship: \texttt{ate.spec} sweeps every non-\texttt{.py} file under
\texttt{src/ate}, and \texttt{[tool.setuptools.package-data]} lists \texttt{*.pdf}, so the frozen bundle and
the wheel carry the same files. \texttt{verify\_bundle()} then refuses a bundle whose PDF is missing or
empty when \texttt{ATE\_\allowbreak{}RELEASE\_\allowbreak{}BUILD=1} is in the environment, which is what the release jobs set;
a build from a plain checkout says plainly that it has no manuals and carries on.

The names come from one list, \texttt{MANUAL\_\allowbreak{}NAMES} in
\texttt{docs/\allowbreak{}manuals/\allowbreak{}build.py}, and \texttt{packaging/\allowbreak{}manuals.py} reads it. A directory with a
\texttt{main.tex} that is not in the list makes the build fail, so a fourth manual cannot be left out
of the artifacts silently. The \texttt{.deb} also copies the PDFs to
\texttt{/usr/share/\allowbreak{}doc/ate}, where Debian expects documentation, and \texttt{ate --help} prints the
path the running install really has: \texttt{\_internal/ate/manuals} inside a frozen bundle,
\texttt{ate/manuals} inside the wheel's package, \texttt{/usr/share/doc/ate} in a \texttt{.deb}. On a tag
the \texttt{release} job attaches all six files --- the three PDFs and the three Markdown manuals ---
to the GitHub Release, each asset named with the version it belongs to, taken from the artifacts this
same run built and smoke-tested.

\section{The CI that runs all of this}

\texttt{.github/\allowbreak{}workflows/\allowbreak{}ci.yml} runs the suite as a single job with a
$2 \times 2$ matrix --- \texttt{ubuntu-latest} and \texttt{windows-latest}, each on CPython 3.13 and
3.14 --- so four job runs, \texttt{fail-fast: false}, with
\texttt{QT\_\allowbreak{}QPA\_\allowbreak{}PLATFORM=offscreen} and \texttt{PYTHONFAULTHANDLER=1} set at job level. Two things in it
are load-bearing and were added because of failures, not taste:

\begin{itemize}
  \item \texttt{timeout-minutes} is 30 on Linux and 40 on Windows. The first real run of this workflow
    hung on \texttt{windows-latest} for an hour on a modal dialog waiting for an Operator who was not
    there. A bounded job that reports a timeout plus a \texttt{faulthandler} stack is a job you can read.
  \item A step checks that \texttt{PySide6-Addons}, \texttt{PySide6-Essentials} and \texttt{shiboken6} are
    the same version and that the Qt the runtime reports matches them. Two Qt builds in one process is how
    the first Windows run died with \texttt{0xc0000139} (STATUS\_ENTRYPOINT\_NOT\_FOUND) the moment
    WebEngine was imported: pip had put a Qt 6.12 core next to a Qt 6.11 WebEngine.
\end{itemize}

\texttt{.github/\allowbreak{}workflows/\allowbreak{}release.yml} runs on a tag or by hand: \texttt{manuals} builds the PDFs
first, \texttt{windows-installer} and \texttt{linux-build} take them as input, and \texttt{linux-smoke}
(matrix: \texttt{deb}, \texttt{tar}, \texttt{wheel}) installs each artifact in a fresh \texttt{debian:stable}
container --- four job definitions, six job runs. On a tag only, a fifth job, \texttt{release}, creates
the GitHub Release for that tag out of the artifacts this same run uploaded. The smoke scripts are
\texttt{packaging/\allowbreak{}linux/\allowbreak{}smoke\_\allowbreak{}linux.py} and
\texttt{packaging/\allowbreak{}windows/\allowbreak{}smoke\_\allowbreak{}installer.py};
\texttt{packaging/\allowbreak{}linux/\allowbreak{}docker\_\allowbreak{}smoke.py} is the driver a human runs.
Each smoke run proves the manuals came with the artifact it installed: \texttt{ate --help} names the
directory, every name on the list is in it, each file starts with \texttt{\%PDF}, and each title carries
the tag's version. The \texttt{.deb} run looks in \texttt{/usr/share/\allowbreak{}doc/ate} too, and the
Windows run checks the bundle after the install and again after the upgrade.
The driver tags a fresh copy of the checkout to build against: \texttt{ATE\_\allowbreak{}SMOKE\_\allowbreak{}TAG}
names the tag, and with nothing named it is this checkout's own \texttt{git describe --tags}. The
staged manuals are then checked against that version rather than supplying it; read from the PDFs,
the expectation would be the thing being checked, and manuals left over from the last release would
make every check in the run agree with itself. A container built as \texttt{1.2.3} while the staged
manuals say something else is a version mismatch, not a packaging failure, and the run refuses it
rather than pass it.

\section{What has actually been run here}

Both gates were run for this manual, on the machine that wrote it, with \texttt{TMPDIR} pointed at a
directory on disk rather than \texttt{/tmp}. \texttt{ATE\_\allowbreak{}BUILD\_\allowbreak{}TESTS=1 pytest -m build} takes a few minutes
because it builds a frozen bundle; \texttt{ATE\_\allowbreak{}DOCKER\allowbreak{}TESTS=1 pytest -m build} takes several, because it builds the tarball
and the \texttt{.deb} and installs each in a fresh \texttt{debian:stable} container. Both pass. The first
attempt at the container gate in this environment died with \texttt{OSError: [Errno 28] No space left on
device} on a small \texttt{/tmp}, which is why the manual says where a gate ran and not merely that it
can. The timings quoted in chapter~\ref{ch:intro} are three separate runs of three
separate commands on one machine, not a comparison.

Packet 36's gates ran here the same way, and the numbers are this machine's. The three PDFs staged
into \texttt{src/\allowbreak{}ate/\allowbreak{}manuals} are just under 0.8 MB. The bundle built with them is 447 files and
489.6 MB; the same bundle with the three files taken out is 444 files and 488.8 MB, so the manuals
cost 0.17\% of a bundle. Packaged from the same bundle: the tarball gains 0.80 MB (199.8 MB to
200.6 MB), the \texttt{.deb} 1.6 MB (148.1 MB to 149.7 MB) --- the \texttt{.deb} pays twice, because the
PDFs ride in the bundle and are copied to \texttt{/usr/share/\allowbreak{}doc/ate} --- and the wheel gains 0.80 MB
(304 KB to 1.11 MB), which is 264\% of it: the wheel ships no Qt, so the manuals are the largest thing
in it.

What has never been run on
this machine is a Windows build, a Windows install, or an installer: the release runner did compile
\texttt{ate.iss} and install, run, upgrade and uninstall the installer at \texttt{v0.1.0-rc2} (run
41198438867), so those paths are proven somewhere --- just not here.

\cando{say which artifact carries which library, and name the test that proves a bundle is complete.}
